\section{扩散模型的完整图景}

\subsection{前向过程与逆过程的对称性}

扩散模型的完整框架可以用以下对称结构概括：

\begin{table}[H]
\centering
\begin{tabular}{lll}
\toprule
 & \textbf{前向过程（加噪）} & \textbf{逆过程（去噪）} \\
\midrule
方向 & $x_0 \to x_1 \to \cdots \to x_T$ & $x_T \to x_{T-1} \to \cdots \to x_0$ \\
分布 & $q(x_t | x_{t-1})$，预定义 & $p_\theta(x_{t-1} | x_t)$，需学习 \\
是否训练 & 否 & 是 \\
终点 & $x_T \approx \mathcal{N}(0, I)$ & $x_0 \sim p_{\text{data}}$ \\
作用 & 训练时产生噪声样本 & 生成时产生数据样本 \\
\bottomrule
\end{tabular}
\caption{前向过程与逆过程的对比}
\end{table}

\subsection{为什么扩散模型有效？}

\begin{knowledgebox}{从 VAE 的视角理解扩散模型}
扩散模型可以被理解为一种特殊的层级 VAE，具有以下关键特性：
\begin{itemize}
    \item 编码器（前向过程）是固定的，不需要学习
    \item 所有层级的维度相同
    \item 每一步的转移都是简单的高斯扰动
    \item 通过足够多的步骤，简单的高斯转移的乘积可以表达非常复杂的分布
\end{itemize}
这种设计在简单性和表达能力之间取得了优雅的平衡。
\end{knowledgebox}

\subsection{后续内容预告}

本讲建立了 DDPM 的概念框架，下一讲（Lecture 03：DDPM Part 2）将继续深入：

\begin{itemize}
    \item 一致性项 $L_{t-1}$ 的详细推导
    \item $q(x_{t-1}|x_t, x_0)$ 的闭式高斯表达式
    \item 从预测均值到\textbf{预测噪声} $\epsilon_\theta$ 的简化
    \item DDPM 的最终训练目标：简单的加权 MSE 损失
    \item 实际训练和采样算法
\end{itemize}

\subsection{本章小结}

扩散模型通过固定前向过程（加噪）和学习逆过程（去噪），在层级 VAE 的框架内实现了强大的生成能力。前向过程的精巧设计使得 ELBO 的多个项自动满足或可以高效计算，将训练聚焦于核心的一致性项。


